\chapter{The Central Risk Book}\label{ch:s2:the-central-risk-book}

Ten desks each hedging their own risk pay ten spreads for trades that would net to little. A \omterm{def:s2:the-central-risk-book:crb}{central risk book} collects the risk first, crosses what
offsets, and hedges only what is left, when it is cheap to do so. On this chapter's synthetic bank, five desks receiving independent client flows in
100 stocks would pay \$335,000 a day to hedge each flow at once. Pooled, the net costs \$270,000, 19.5\% less. Kept and traded out at 15\% a day, with
the market risk hedged in index futures, it costs \$45,000 a day. The book then carries residual risk with a one-day 99\% value at risk of \$859,000.
The build is \texttt{firm.crbook}.

\section{Pooling and netting}

\begin{definition}[Central risk book]\label{def:s2:the-central-risk-book:crb}
A \emph{central risk book}\index{central risk book} is a trading book that takes on the risk that several desks acquire from clients, nets it, and
manages the net: crossing offsetting positions internally, hedging common factors with liquid instruments and trading out the rest over time within
risk limits.
\end{definition}

\begin{definition}[Risk pooling]\label{def:s2:the-central-risk-book:pooling}
\emph{Risk pooling}\index{risk pooling} is combining positions from several sources before hedging, so that offsetting exposures cancel and only the
net is traded.
\end{definition}

\texttt{firm.crbook} gives five desks client flows in 100 stocks for 500 days (\cref{lst:s2:the-central-risk-book:book}). Each desk's flow in each name
has a standard deviation of \$0.5 million a day, and 10\% of the variance is common to all desks, as when clients herd. Stock returns follow a
three-factor model, with a market factor of 1\% daily volatility, two style factors and 1.5\% of idiosyncratic volatility. Trading costs use Book~7's
model (chapter~27): a half-spread of 3 basis points and square-root impact with a coefficient of 0.7, on daily volumes between \$20 and \$200 million.
Almgren, Thum, Hauptmann and Li, from almost 700,000 orders, found temporary impact closer to a 3/5 power of the trade rate than a square root; the
exponent matters less here than the fact that impact grows faster than size.

Hedged desk by desk, each flow executed on its own as if the desk were alone in the market, the hedges cost \$335,000 a day. Pooled first, 47.0\% of the
gross flow crosses internally and the net costs \$270,000, a saving of 19.5\%. The saving is smaller than the share crossed because the impact per
dollar is lower on small trades: what nets away is disproportionately the cheap part. Butz and Oomen described the choice every dealer makes between
internalising a client's risk, warehousing it in anticipation of offsetting flow, and externalising it by hedging at once; a \omterm{def:s2:the-central-risk-book:crb}{central risk book} makes
that choice for the whole bank.

\section{Internal crossing}

The saving depends on how independent the desks' flows are (\cref{fig:s2:the-central-risk-book:desks}). With fully independent flows, pooling ten desks
saves 48.0\% of the cost. With 10\% of the flow common to all desks, the saving peaks at 19.7\% with six desks and falls to 17.1\% with ten: the common
part does not net, and with more desks it becomes a larger share of the net. A bank's desks see the same clients' moods, and the case for pooling must
be made on their actual correlation.

\begin{omfigure}
\begin{tikzpicture}
\begin{axis}[width=11cm,height=5.2cm,xlabel={\small number of desks pooled},ylabel={\small hedging cost saved (\%)},
  tick label style={font=\footnotesize},xmin=1,xmax=10,ymin=0,ymax=50,grid=major,grid style={gray!25},table/col sep=comma,
  legend style={font=\scriptsize,at={(0.97,0.03)},anchor=south east}]
\addplot[omThm,thick,mark=*,mark size=1.2pt] table[x=desks,y=independent]{figdata/strategies-2/25-the-central-risk-book/desks.csv};
\addplot[omDef,thick,mark=square*,mark size=1.2pt] table[x=desks,y=common]{figdata/strategies-2/25-the-central-risk-book/desks.csv};
\legend{independent flows,10\% of flow variance common}
\end{axis}
\end{tikzpicture}

\omcaption{Hedging cost saved by pooling desks' flows before hedging them at once, against the number of desks, for independent flows and for flows
with a common component. Data: \texttt{s2\_crbook.by\_desks}.}\label{fig:s2:the-central-risk-book:desks}
\end{omfigure}

\section{Optimised hedging}

\begin{definition}[Residual risk hedge]\label{def:s2:the-central-risk-book:residual}
A \emph{residual risk hedge}\index{residual risk hedge} hedges a pooled book's exposure to common factors at once with liquid instruments (index futures,
factor baskets) and leaves the remaining, mostly idiosyncratic, positions to be traded out gradually or offset by later flows.
\end{definition}

The central book need not hedge the net at once. Each day it takes the clients' net flow into its inventory, sells a fixed fraction of the inventory
in the stocks, and holds an index future against the beta-weighted remainder. Garleanu and Pedersen showed that when trading is costly the optimal
policy trades partially toward its target each period; here the target is a flat book, and the fraction sets the trade-off. Tomorrow's flows offset
part of what is kept, and what is traded out is traded in smaller pieces:

\begin{center}
\small
\begin{tabular}{@{}lccc@{}}
\toprule
share traded out a day & cost (\$ thousands a day) & P\&L sd (\$ thousands) & sd without futures\\
\midrule
100\% (pooled, at once) & 269.7 & 0 & 0\\
50\% & 123.1 & 128 & 157\\
30\% & 78.4 & 222 & 268\\
20\% & 56.5 & 304 & 367\\
15\% & 45.3 & 369 & 448\\
10\% & 33.4 & 472 & 579\\
5\% & 20.3 & 684 & 842\\
\bottomrule
\end{tabular}
\end{center}

The futures overlay costs about \$500 a day and removes about a fifth of the risk: the net positions are spread across names, so most of what is carried
is idiosyncratic. Which point on the frontier (\cref{fig:s2:the-central-risk-book:frontier}) is right depends on the price of risk. With a one-day 99\%
value at risk limit of \$1 million, the cheapest rate within the limit is 15\%: \$45,300 a day and a value at risk of \$859,000. With a limit of
\$500,000, it is 35\%: \$89,300 a day. Against hedging desk by desk, the first saves 86.5\% of the cost.

\begin{omfigure}
\begin{tikzpicture}
\begin{axis}[width=11cm,height=5.4cm,xlabel={\small standard deviation of the daily P\&L carried (\$ thousands)},
  ylabel={\small hedging cost (\$ thousands a day)},tick label style={font=\footnotesize},xmin=0,xmax=900,ymin=0,ymax=290,grid=major,
  grid style={gray!25},table/col sep=comma,legend style={font=\scriptsize,at={(0.97,0.97)},anchor=north east}]
\addplot[omThm,thick,mark=*,mark size=1pt] table[x=sd,y=cost]{figdata/strategies-2/25-the-central-risk-book/frontier.csv};
\addplot[gray,thick,dashed,mark=square*,mark size=1pt] table[x=sd_nofut,y=cost_nofut]{figdata/strategies-2/25-the-central-risk-book/frontier.csv};
\addplot[omDef,thick,dotted] coordinates {(429.9,0) (429.9,290)};
\legend{with index-future overlay,stocks only,value-at-risk limit \$1m}
\end{axis}
\end{tikzpicture}

\omcaption{The central book's cost-risk frontier as the share of inventory traded out each day runs from 5\% to 100\%, with and without an index-future
overlay; the dotted line is the standard deviation at which the one-day 99\% value at risk reaches \$1 million. Data:
\texttt{s2\_crbook.frontier}.}\label{fig:s2:the-central-risk-book:frontier}
\end{omfigure}

\section{Governance}

A \omterm{def:s2:the-central-risk-book:crb}{central risk book} changes who owns risk and who pays for hedging. The desks that bring the flows must be charged a transfer price for the risk they
hand over, and credited for the netting their flows provide. If the price is too low, desks send risk they should have priced out of their clients; if
it is too high, they hedge on their own and the netting is lost. The book's own limits (value at risk, factor exposures, concentration, time to
liquidate) need independent risk management, and the crossing itself needs rules. Clients' trades must be executed on the terms promised to them, and
information from one desk's clients must not reach another desk's trading. Book~7 (chapter~27) treats the internal crossing of a firm's own strategies;
the same questions of fair allocation arise between desks.

\section{Strategy files}

\begin{strategyfile}{Central netting across desks}\label{strat:s2:the-central-risk-book:netting}
\sfield{Who pays you, and why} The market's spread and impact, not paid on offsetting flows.
\sfield{Instruments and venues} Desks' positions; the internal crossing engine.
\sfield{Signal} Offsetting exposures across desks.
\sfield{Sizing and execution} Net before hedging, with transfer prices to the desks.
\sfield{Costs} Systems and governance.
\sfield{How it dies} Correlated flows that do not net; desks that route around it.
\sfield{Horizon, capacity, infrastructure} Intraday; a firm-wide position feed.
\sfield{Backtest honestly} Desks' actual flows and their correlation, not independence.
\sfield{Sources} Butz and Oomen (2019); this chapter: 19.5\% of hedging cost saved by pooling five desks.
\end{strategyfile}

\begin{strategyfile}{Factor-hedged residual risk}\label{strat:s2:the-central-risk-book:factor}
\sfield{Who pays you, and why} The cheapness of index futures against single stocks.
\sfield{Instruments and venues} Index futures, factor baskets, sector ETFs.
\sfield{Signal} The pooled book's factor exposures from a risk model.
\sfield{Sizing and execution} Hedge factors at once; leave idiosyncratic risk.
\sfield{Costs} Futures roll and basis.
\sfield{How it dies} A risk model that misses a factor.
\sfield{Horizon, capacity, infrastructure} Daily; a factor risk model (Book~7, chapter~24).
\sfield{Backtest honestly} Out-of-sample factor exposures.
\sfield{Sources} This chapter: a fifth of the carried risk removed for \$500 a day.
\end{strategyfile}

\begin{strategyfile}{Patient hedging of pooled risk}\label{strat:s2:the-central-risk-book:patient}
\sfield{Who pays you, and why} Impact saved by trading smaller and letting later flows offset.
\sfield{Instruments and venues} The stocks, traded over days.
\sfield{Signal} Inventory against the cost-risk frontier.
\sfield{Sizing and execution} Trade a fixed share of the inventory a day, set by the risk limit.
\sfield{Costs} Residual risk.
\sfield{How it dies} One-way flows; volatility spikes that break the limit.
\sfield{Horizon, capacity, infrastructure} Days; risk limits and capital.
\sfield{Backtest honestly} Flows and volatility of stressed periods, not average ones.
\sfield{Sources} Garleanu and Pedersen (2013); this chapter: \$45,300 a day at a \$1 million value at risk, against \$335,000 desk by desk.
\end{strategyfile}

\section{Tutorial: hedge what is left}\label{tut:s2:the-central-risk-book:tut}

\begin{tutorial}
\textbf{Goal.} Simulate desks' client flows, hedge them desk by desk, pooled and patiently with a futures overlay, and choose the hedging rate from a
risk limit. \textbf{End state:} the table and two figures.
\begin{tutsteps}
\item \textbf{The book}.
\omcode{code/firm/crbook/firm_crbook.py}{62}{98}{Hedging cost, desk by desk, pooled, and the patient central book.}\label{lst:s2:the-central-risk-book:book}
\item \textbf{The frontier and the limit}.
\omcode{code/strategies-2/25-the-central-risk-book/python/s2_crbook.py}{42}{59}{Cost and risk by rate; the cheapest rate within a value-at-risk limit.}\label{lst:s2:the-central-risk-book:frontier}
\item \textbf{Run} \texttt{pooling()}, \texttt{frontier()}, \texttt{pick()}, \texttt{by\_desks()} and \texttt{fig\_crbook.py}.
\end{tutsteps}
\textbf{What to change next.} Make one desk's flow a hedge of another's (negative correlation); trade out names at rates set by their liquidity; add a
transfer price and see which desks would route around the book.
\end{tutorial}

\section{Build: central risk book}\label{bld:s2:the-central-risk-book:crbook}

\begin{build}
\bfield{Purpose} Desks' flows, pooling, internal crossing, a futures overlay and patient hedging, costed with Book 7's model.
\bfield{Interface} \texttt{CRBConfig(\dots)}, \texttt{simulate\_desks(cfg)}, \texttt{hedge\_cost(trade, sim, cfg)}, \texttt{desk\_by\_desk},
\texttt{pooled}, \texttt{central(sim, cfg, rate, futures)}.
\bfield{Rules} Desks' hedges executed as if alone; the central book trades a fixed share of inventory a day; futures against beta-weighted
inventory.
\bfield{Acceptance tests} \texttt{code/firm/crbook/tests/}: offsetting desks cost nothing pooled; trading everything at once is pooling and carries
nothing; futures remove market risk; the cost by hand.
\bfield{Stretch} Name-by-name rates; transfer pricing; stressed flows.
\end{build}

\begin{omsources}
\item N.~Garleanu and L.~H.~Pedersen, ``Dynamic trading with predictable returns and transaction costs'', \emph{Journal of Finance} 68(6), 2013.
\item R.~Almgren, C.~Thum, E.~Hauptmann and H.~Li, ``Direct estimation of equity market impact'', \emph{Risk}, 2005.
\item M.~Butz and R.~Oomen, ``Internalisation by electronic FX spot dealers'', \emph{Quantitative Finance} 19(1), 2019.
\end{omsources}

\section{Exercises}

\begin{exercise}[$\star$]\label{exo:s2:the-central-risk-book:1}
Two desks must buy \$3 million and sell \$2 million of the same stock. How much crosses internally, and what share of the gross?
\end{exercise}

\begin{exercise}[$\star$]\label{exo:s2:the-central-risk-book:2}
The carried P\&L has a standard deviation of \$369,000 a day. What is the one-day 99\% value at risk under normality?
\end{exercise}

\begin{exercise}[$\star$]\label{exo:s2:the-central-risk-book:3}
With a half-spread of 3 basis points, impact $0.7\sigma\sqrt{q/V}$, $\sigma = 2\%$ and $V = \$50$ million, what does hedging \$1 million cost?
\end{exercise}

\begin{exercise}[$\star\star$]\label{exo:s2:the-central-risk-book:4}
Why does pooling save less than the share of flow that crosses?
\end{exercise}

\begin{exercise}[$\star\star$]\label{exo:s2:the-central-risk-book:5}
Why does trading out only part of the inventory each day save so much cost?
\end{exercise}

\begin{exercise}[$\star\star$]\label{exo:s2:the-central-risk-book:6}
Why does the pooling saving fall beyond six desks when flows have a common part?
\end{exercise}

\begin{exercise}[$\star\star\star$]\label{exo:s2:the-central-risk-book:7}
\emph{Coding.} Run \texttt{pick(5e5)}. Which rate does a \$500,000 value-at-risk limit choose, and what does it cost?
\end{exercise}

\begin{exercise}[$\star\star\star$]\label{exo:s2:the-central-risk-book:8}
\emph{Find the flaw.} ``The central book cut hedging costs by 86\%; we should trade out 5\% a day and cut them by 94\%.''
\end{exercise}

\section{Problem: Hedge What Is Left}

\begin{problem}[{Weekend problem --- the central risk book}]\label{pb:s2:the-central-risk-book:1}
The chapter's synthetic desks and the public record.

\medskip\noindent\textbf{Part I --- Pooling.}
\begin{enumerate}
\item Define a \omterm{def:s2:the-central-risk-book:crb}{central risk book} and \omterm{def:s2:the-central-risk-book:pooling}{risk pooling}.
\item Describe the desks, flows, returns and costs.
\item Give the costs desk by desk and pooled, and the share crossed.
\item What choice did Butz and Oomen describe?
\end{enumerate}

\medskip\noindent\textbf{Part II --- Crossing.}
\begin{enumerate}[resume]
\item How does the saving depend on the number of desks?
\item Why does the common part matter?
\item What did Almgren and co-authors find about impact?
\item Why is the saving smaller than the share crossed?
\end{enumerate}

\medskip\noindent\textbf{Part III --- Optimised hedging.}
\begin{enumerate}[resume]
\item Define a \omterm{def:s2:the-central-risk-book:residual}{residual risk hedge}.
\item Describe the patient policy and give the frontier.
\item What does the futures overlay do?
\item Which rate does a \$1 million limit choose?
\end{enumerate}

\medskip\noindent\textbf{Part IV --- The verdict.}
\begin{enumerate}[resume]
\item State the \emph{named result}: the hedging cost saved by pooling and the residual risk accepted.
\item Which assumption flatters the desk-by-desk comparison?
\item What should the transfer price to desks do?
\item What can go wrong with a patient book in a one-way market?
\item What must governance keep apart?
\item Which strategy file depends on a risk model?
\item How does this chapter relate to chapter~24's internalisation?
\item In one sentence: what does a \omterm{def:s2:the-central-risk-book:crb}{central risk book} sell to the bank?
\end{enumerate}
\end{problem}

\section{Interview questions}

\begin{interviewq}[$\star$ \iqroles{trader}]\label{iq:s2:the-central-risk-book:1}
What is a \omterm{def:s2:the-central-risk-book:crb}{central risk book}, and why would a bank run one?
\end{interviewq}

\begin{interviewq}[$\star\star$ \iqroles{researcher}]\label{iq:s2:the-central-risk-book:2}
How would you choose how fast to trade out a pooled inventory?
\end{interviewq}

\begin{interviewq}[$\star\star$ \iqroles{trader}]\label{iq:s2:the-central-risk-book:3}
Two desks send you opposite risk in the same stock at the same time. How do you price the cross to each?
\end{interviewq}

\begin{interviewq}[$\star\star$ \iqroles{risk}]\label{iq:s2:the-central-risk-book:4}
What limits would you put on a \omterm{def:s2:the-central-risk-book:crb}{central risk book}?
\end{interviewq}

\begin{interviewq}[$\star\star$ \iqroles{developer}]\label{iq:s2:the-central-risk-book:5}
Design the firm-wide position feed a \omterm{def:s2:the-central-risk-book:crb}{central risk book} needs.
\end{interviewq}

\begin{interviewq}[$\star\star\star$ \iqroles{researcher}]\label{iq:s2:the-central-risk-book:6}
$D$ desks have independent flows of standard deviation $s$ in a name. Show that pooling cuts the expected linear cost by a factor $\sqrt D$ and
the expected impact cost, proportional to $|q|^{3/2}$, by a factor $D^{1/4}$.
\end{interviewq}
